\begin{description}
\item[a)] As it was suggested in the hint the key is a comparison of the
sunlight scattered from a $10\,\mathrm{cm} \times 10\,\mathrm{cm}$ aluminium
plate and the Moon to the observer (both are nearly at the local zenith).

Let the distance of the Moon and the satellite from the observing site be
$d_{\mathrm{Moon}}$, and $d_{\mathrm{MASAT}}$, respectively. Under the
mentioned "ideal" circumstances, when both the Moon and the satellite is near
the local zenith: $d_{\mathrm{MASAT}} = 900$\,km, while
\begin{equation*}
d_{\mathrm{Moon}} = s_{\mathrm{Moon}} - R_{\mathrm{Earth}}
  - R_{\mathrm{Moon}}, \tag{14.1}
\end{equation*}
\hfill(2 p)\\
where the average distance between the centre of Earth and Moon
$s_{\mathrm{Moon}}$ is included in the \emph{Table of constants}.

Let the intensity coming from the Sun and scattered back to the direction of
the observer be $I$, this is considered practically the same at the solar
distance of the Moon and the satellite, but necessarily multiplied by the
albedo of the given surface quality (the lunar albedo $a_{\mathrm{Moon}}$ is
given in the Constant table, while $a_{\mathrm{MASAT}}$ of the satellite is
in the text).

Finally, for using the magnitude formula, to find the observable magnitude of
the satellite, we have to consider the resulted flux of the scattered light
at the observer, which is directly proportional to the intensity, the albedo
and the scattering area, but is inversely proportional to the distance on the
square:
\begin{equation*}
I_{\mathrm{scat}} \approx
  \frac{a_{\mathrm{obj}}\,I\,A_{\mathrm{obj}}}{d_{\mathrm{obj}}^2}
\tag{14.2}
\end{equation*}
\hfill(3 p)

Thus, the magnitude formula gives:
\begin{equation*}
m_{\mathrm{MASAT}} - m_{\mathrm{Moon}}
  = -2.5\log\frac{I_{\mathrm{MASAT}}}{I_{\mathrm{Moon}}}
  = -2.5\log\frac{\dfrac{a_{\mathrm{MASAT}}\,I\,A_{\mathrm{MASAT}}}
      {d_{\mathrm{MASAT}}^2}}
    {\dfrac{a_{\mathrm{Moon}}\,I\,A_{\mathrm{Moon}}}{d_{\mathrm{Moon}}^2}}
\tag{14.3}
\end{equation*}
\hfill(3 p)

With numerical values:
\[ m_{\mathrm{MASAT}} = 9.72^{\mathrm{m}} \]
\hfill(2 p)

\item[b)] A moving point-like object will generally cause a line on the CCD
image, resulting less photoelectrons in a given pixel, than a similarly
bright unmoving object. Since the MASAT--1 satellite is so small, it is not
resolved by the telescope -- it will cause a seeing-sized spot at the focus
of a telescope, like any other stars on the sky, but since the text of the
problem described: we should omit the seeing. Thus, the image of any point
source will be formed by the diffraction -- so we shall use the approximate
formula of the Airy disk.

By the diffraction-limited imaging, the point-like sources illuminate a spot
in the focal plane having the angular diameter of the so-called Airy disc
determined by the formula:
\begin{equation*}
\varphi_{\mathrm{Airy}} = 2.44\,\frac{\lambda}{D}, \tag{14.4}
\end{equation*}
\hfill(2 p)\\
while its linear size at the focus of $f$ focal length telescope is:
\begin{equation*}
d_{\mathrm{Airy}} = f\varphi_{\mathrm{Airy}} \tag{14.5}
\end{equation*}
\hfill(2 p)

For the telescope of Baja Observatory, described in the text, at a typical
visual wavelength of $0.55\,\mu\mathrm{m}$:
\[ d_{\mathrm{Airy}} = 2.44\,\frac{f\lambda}{D}
   = 2.44 \times \frac{4200\,\mathrm{mm} \times 0.55\,\mu\mathrm{m}}
     {500\,\mathrm{mm}}
   = 11.27\,\mu\mathrm{m} \]
\hfill(2 p)

This is very close to the pixel size of the applied CCD. However, this
number is not necessary for our calculation, since both the stars and the
MASAT--1 are point-like sources, illuminating the same number of pixels, thus
the number of pixels will drop out from all comparative calculations. So
does the quantum efficiency value of the applied CCD.

For deriving the effect of the moving object, let us first calculate the
scale of the image at the focus of the 50\,cm RC telescope of Baja ($N$ is
the f-number of the telescope, $N = f/D$):
\[ S = \frac{206\,265''}{DN}
     = \frac{206\,265''}{500\,\mathrm{mm} \times 8.4}
     = 49\,''\,\mathrm{mm^{-1}} = 0.049\,''\,\mu\mathrm{m}^{-1} \]
\hfill(3 p)

This means, that one pixel of the attached CCD describes
$P = 0.049\,''\,\mu\mathrm{m}^{-1} \times 9\,\mu\mathrm{m} = 0.44''$
side length area on the sky. \hfill(3 p)

Now let us estimate the relative flux, resulted by a star having the limiting
magnitude of the telescope-CCD system ($19.5^{\mathrm{m}}$), by the magnitude
formula:
\begin{equation*}
19.5 = -2.5\log\Phi_{19.5} \rightarrow
  \Phi_{19.5} = 10^{-7.8}\,\mathrm{s^{-1}} \tag{14.6}
\end{equation*}
\hfill(2 p)

And similarly, for the MASAT--1 brightness -- observed with the same system,
under the same conditions, if it were an unmoving object:
\begin{equation*}
9.72 = -2.5\log\Phi_{\mathrm{MASAT}} \rightarrow
  \Phi_{\mathrm{MASAT}} = 10^{-3.9}\,\mathrm{s^{-1}} \tag{14.7}
\end{equation*}
\hfill(2 p)

Next, we have to calculate the angular speed of MASAT--1 at 900\,km altitude
above the observing site, in order to derive the fluxes falling onto one
pixel of the CCD on the telescope.

The velocity of the motion on a circular orbit with orbital radius $r$:
\begin{equation*}
v_r = \sqrt{\frac{GM_{\mathrm{Earth}}}{r}}, \tag{14.8}
\end{equation*}
\hfill(2 p)\\
where $r = h + R_{\mathrm{Earth}}$. \hfill(2 p)

With numerical values:
\[ v_r = 7404\,\mathrm{m\,s^{-1}} \]
\hfill(1 p)

Hence, the visible angular speeds (aside from the Earth rotation):
\begin{equation*}
\omega_r = \frac{v_r}{r} \tag{14.9}
\end{equation*}
\hfill(2 p)
% sic: the source prints omega_r = v_r / r; for the angular speed seen from
% the observing site directly below, the distance should be the altitude h.
With numerical values:
\[ \omega_r = 0.47\,\mathrm{deg\,s^{-1}} = 1696.8\,''\,\mathrm{s^{-1}}, \]
\hfill(1 p)\\
when the satellite is near the local zenith.

For deriving the relative speed, which will be really exposed onto the CCD,
we have to calculate the angular speed of the observing site caused by the
terrestrial rotation. At the latitude of $\varphi$ it is given as:
\begin{equation*}
\omega_{\mathrm{rot},\varphi} = \omega_{\mathrm{rot,eq}}\cos\varphi
\tag{14.10}
\end{equation*}
\hfill(2 p)

For Baja its numerical value:
\[ \omega_{\mathrm{rot,B}} = 10.4\,''\,\mathrm{s^{-1}} \]
\hfill(1 p)

Its component parallel to the direction of the satellite's orbital motion:
\begin{equation*}
\omega_{\mathrm{rot,B},i} = \omega_{\mathrm{rot,B}}\cos i_{\mathrm{MASAT}}
\tag{14.11}
\end{equation*}
\hfill(2 p)

For Baja its numerical value:
\[ \omega_{\mathrm{rot,B},i} = 3.6\,''\,\mathrm{s^{-1}} \]
\hfill(1 p)

Finally, the relative angular speed of the nanosatellite above the telescope:
\begin{equation*}
\Delta\omega = \omega_a - \omega_{\mathrm{rot,B},i} \tag{14.12}
\end{equation*}
\hfill(2 p)

With numerical values:
\[ \Delta\omega = 1693.3\,''\,\mathrm{s^{-1}} \]
\hfill(1 p)

% FIGURE NEEDED: sketch of the satellite (orbital speed v_a) passing over the
% CCD pixel grid (9 um pixels), with the Airy-disc sized image of the
% nanosatellite moving with angular speed Delta-omega across one pixel and
% Airy-disc sized star images at fixed positions; 2019_th_sol.pdf page 26.

Since the satellite image is moving with a given $\Delta\omega$ angular
speed, this means that the flux coming from the satellite is falling onto a
given pixel only for a very limited time (independently of the exposure time
of the image, but acting similarly, than that). Considering our very simple
approximation we get an approximate "illumination time" for one pixel falling
onto the line of the orbit:
\begin{equation*}
\tau = \frac{P}{\Delta\omega} \tag{14.13}
\end{equation*}
\hfill(2 p)

With numerical values:
\[ \tau = \frac{0.44''}{1693.3\,''\,\mathrm{s^{-1}}}
        = 2.61 \times 10^{-4}\,\mathrm{s} \]
\hfill(1 p)

For finishing our estimation -- to get the magnitude level of any pixel on
the observable trace of this nanosatellite, let us use the magnitude formula
for them and one non-moving light source at the detection limit. For doing
this we have to multiply the relative fluxes $\Phi_{\mathrm{MASAT}}$ and
$\Phi_{19.5}$ with the respective illumination times. (For a
$19.5^{\mathrm{m}}$ stars we should use the 2\,min-long exposure time
mentioned in the text, which is needed to collect enough photoelectrons for
the safe detection of such a faint star.)
\begin{equation*}
m_{\mathrm{MASAT,trace}} - 19.5^{\mathrm{m}}
  = -2.5\log\frac{\tau\Phi_{\mathrm{MASAT}}}
    {\tau_{\mathrm{exp}}\Phi_{19.5}} \tag{14.14}
\end{equation*}
\hfill(2 p)

With numerical values:
\[ m_{\mathrm{MASAT,trace}} = 23.9^{\mathrm{m}} \]
\hfill(2 p)

So, the trace of the moving $23.9^{\mathrm{m}}$ magnitude nanosatellite looks
so faint on a (2\,min-long exposure) CCD image, as a $\sim 24^{\mathrm{m}}$
stars would be seen, but they fall surely below the detection limit for the
described telescope. Thus, the answer is definitely

\textbf{NO}, WE COULD NOT RECORD THE MASAT--1 \hfill(2 p)\\
using that telescope.

\item[c)] If we consider a real atmosphere, due to the effect of the
atmospheric turbulences (the so-called seeing) during the 2 minutes-long
exposure the images of the stars will never be described by the ideal Airy
disk, but expanded to a much larger "seeing spot", which can be approximated
by a Gaussian profile. As mentioned in the text: in Hungary, its diameter at
half-maxima is typically $3.5''$. This means, that the flux coming from a
$19.5^{\mathrm{m}}$ star will be distributed over about a
\[ \frac{\mathrm{FWHM}}{P} \times \frac{\mathrm{FWHM}}{P}
   = \frac{3.5''}{0.44''} \times \frac{3.5''}{0.44''} \approx 8 \times 8 \]
\hfill(1 p)\\
pixels area. Thus, now the $F_{19.5}$ flux will be distributed amongst
$4 \cdot 4\pi \approx 50$ pixels (for the sake of simplicity:
homogeneously), i.e.\ one pixel gets:
\begin{equation*}
\Phi_{19.5} = \frac{10^{-7.8}}{50}\,\mathrm{s^{-1}}
  \approx 10^{-9.5}\,\mathrm{s^{-1}} \tag{14.15}
\end{equation*}
\hfill(1 p)

As we have seen above, the satellite image is moving so fast along the focal
plane (0.26\,ms per a pixel), that during this short time, the atmosphere can
be considered to be still. Thus, the image of MASAT--1 can further be
approximated with the Airy-disc size also in a real atmosphere.

Let us correct our last magnitude formula for this fact:
\begin{equation*}
m_{\mathrm{MASAT,trace}} - 19.5^{\mathrm{m}}
  = -2.5\log\frac{\tau\Phi_{\mathrm{MASAT}}}
    {\tau_{\mathrm{exp}}\dfrac{\Phi_{19.5}}{50}} \tag{14.16}
\end{equation*}
\hfill(2 p)

With numerical values:
\[ m_{\mathrm{MASAT,trace}} = 19.6^{\mathrm{m}} \]
\hfill(2 p)

Under real conditions the situation is much better: the trace of the small
nanosatellite is near the detectability! Thus, owing to further secondary
effects (pixel sensitivity inhomogenities, temporarily changing transparency
of the atmosphere, etc.), the trace of the cubesat can easily be seen on the
CCD of the 50\,cm reflector. So in real cases, the answer can be YES.

\textbf{YES}, WE COULD HAVE RECORDED THE MASAT--1. \hfill(2 p)
\end{description}
